% TEMPORARY working report -- not the paper.
%
% Compile:  latexmk -pdf reports/four_way_temp.tex
% Tables:   python scripts/qa_make_four_way_tex.py
%
% Every number is generated by qa_make_four_way_tex.py from the per-seed JSON
% that scripts/qa_four_way.py emits. Nothing here is transcribed by hand.
\documentclass[11pt,a4paper]{article}

\usepackage[margin=1in]{geometry}
\usepackage{booktabs}
\usepackage{amsmath,amssymb}
\usepackage{adjustbox}
\usepackage{caption}
\usepackage[colorlinks=true,linkcolor=black,urlcolor=blue]{hyperref}
\captionsetup{font=small,labelfont=bf,skip=4pt}
\setlength{\tabcolsep}{5pt}
\setlength{\parskip}{2pt}

\newcommand{\Etheta}{E_{\Theta}}
\newcommand{\util}{\mathcal{U}}

\input{four_way_tables}

\title{\vspace{-1.5em}\textbf{Four Feature Sets, One Cost Model}\\[0.2em]
       \large Baseline, Alpha158 seed, Arm A and Arm B on CSI~300, \WindowFW}
\author{QuantaAlpha --- \emph{temporary working report}}
\date{$\Theta$ hash \texttt{\ThetaHashFW} \quad construction \ConstructionFW\
      ($\Sigma$=\CovFW, $\lambda$=\LamFW, turnover cap \CapFW, max $w$ \MaxWFW)}

\begin{document}
\maketitle
\vspace{-1em}

\section{What is being compared}

Four feature sets are priced by one evaluation operator $\Etheta$, one portfolio
map $g$ and one cost model, so the only thing that differs between rows is the
features themselves:

\begin{itemize}\itemsep1pt\parskip0pt
  \item \textbf{baseline} --- the four hand-written base features, no mining at
        all. The floor: mining has to clear this to have contributed anything.
  \item \textbf{\texttt{alpha158\_20}} --- Qlib's own 20-factor Alpha158 subset,
        evaluated through Qlib's expression engine rather than translated into
        the mining DSL, so it is the reference rather than a paraphrase of it.
        It owes nothing to this pipeline.
  \item \textbf{Arm A} --- the RankIC arm's mined factors (the published
        objective).
  \item \textbf{Arm B} --- the net-of-cost arm's mined factors.
\end{itemize}

Base features are excluded from the three non-baseline rows. They are common to
all of them, so including them would measure what \emph{they} contribute rather
than what mining did. The consequence is that these are not the books anyone
would deploy; they are the comparison that isolates the feature sets.

Costs are the full model: commission $\kappa_0$, volatility-scaled slippage
$\kappa_1$, super-linear impact $\kappa_2$ and borrow. This is materially
harsher than the flat commission the published numbers use, and the gap between
the two is the point of the exercise rather than a discrepancy to reconcile.

\paragraph{The construction was chosen on validation, not here.}
$g$'s parameters --- risk model, risk aversion, turnover budget, position cap
--- were selected by sweeping 54 configurations on the \textbf{2021 validation
split} using Arm~B's mined factors, and the winner was priced on
\WindowFW\ exactly once afterwards. Selecting a construction by its performance
on the window then reported would consume the split the design exists to
protect. Two results from that sweep bear on how the table below reads: a
$k$-factor statistical risk model was \emph{beaten} by the diagonal one in 28 of
36 matched pairs, and net IR rises monotonically with the turnover budget
($+0.2619$ at $0.004$/day, $+0.2963$ at $0.008$, $+0.3344$ at $0.02$).

\section{Headline}

\begin{table}[h]\centering
\caption{Full period, mean across combiner seeds $(42,1,7)$. Mined factors only.}
\adjustbox{max width=\textwidth}{\HeadlineTable}
\end{table}

Three things are worth stating plainly.

\textbf{Both arms beat both references.} Arm~B returns $-0.60\%$ against the
baseline's $-1.66\%$ and the seed library's $-2.26\%$; Arm~A returns $-1.35\%$.
Mining contributes something that neither the hand-written features nor a
standard factor set supplies.

\textbf{The seed library is the worst row.} \texttt{alpha158\_20} underperforms
doing nothing at all once slippage and impact are charged. It is not a weak
result for Alpha158 so much as a statement about what a generic factor set costs
to trade: its Rank IC ($+0.0190$) is indistinguishable from the baseline's
($+0.0204$) while it pays more to run.

\textbf{Every row here loses money after costs}, the best returning $-0.60\%$.
That is a statement about \emph{this} construction rather than about the
factors: \S\ref{sec:matrix} shows the same Arm~B library returning $+0.78\%$
after the same full costs under top-$k$, which is the one profitable
configuration found anywhere in this report. Read the headline table as the
mean-variance column of that grid, not as the last word on what the libraries
can do.

\section{Is the ordering readable?}

\begin{table}[h]\centering
\caption{Pairwise net-ARR gaps against the largest combiner-seed standard
deviation in either configuration. A gap under $2\times$ the noise cannot
support a ranking.}
\adjustbox{max width=\textwidth}{\ResolveTable}
\end{table}

All six comparisons resolve, including \textbf{Arm~B over Arm~A at
$4.7\times$}. This is a change of status rather than of magnitude: under the
previous construction (max $w=0.05$, $\lambda=10$) the same comparison was
$+0.60$pp at $1.3\times$ noise and was \emph{not} reportable. Tightening the
position cap cut the seed dispersion from $0.45$pp to $0.16$pp, and it is that
--- not a larger effect --- which made the difference legible.

One caveat that the table cannot carry: run-to-run \emph{generation} noise is
far larger than combiner-seed noise. An unchanged arm has moved $4.05$pp of ARR
between mining runs, roughly $9\times$ the seed dispersion. These arms were
mined once each, so the resolvability above is conditional on the libraries;
it is not a claim that re-mining Arm~B would reproduce the gap.

\section{Year by year}

\begin{table}[h]\centering
\caption{Rank IC by calendar year.}
\adjustbox{max width=\textwidth}{\YearRankIC}
\end{table}

\begin{table}[h]\centering
\caption{Net ARR by calendar year, after full costs.}
\adjustbox{max width=\textwidth}{\YearNetARR}
\end{table}

\begin{table}[h]\centering
\caption{Net IR by calendar year.}
\adjustbox{max width=\textwidth}{\YearNetIR}
\end{table}

The full-period figures hide a regime break that matters more than the ordering.
\textbf{Every configuration is profitable in 2022--2023 and loses money in
2024--2025.} Arm~B earns $+2.98\%$ and $+4.56\%$, then $-7.49\%$ and $-3.96\%$.
The factors were mined on data ending in 2021 and then frozen, and the
$-0.60\%$ headline is an average across a regime they understood and one they
did not.

Arm~B's edge also decays more slowly than Arm~A's. In 2024 its Rank IC is
$+0.0404$ against Arm~A's $+0.0281$, and in 2025 $+0.0322$ against $+0.0244$ ---
gaps of $4.1\times$ and $9.8\times$ the seed noise respectively. The
net-of-cost objective is not merely better on average; it is better later, which
is the property that matters if the library is going to be held rather than
re-mined continuously.

\section{Cost model versus construction}
\label{sec:matrix}

The two tables this project has been reporting each hold one cell of a
$2\times2$, and they differ in \emph{both} axes at once: the flat-fee backtest
runs top-$k$ dropout, while $\Etheta$ runs mean-variance under the full cost
model. The distance between them therefore mixes the cost model with the
construction. Filling the grid separates them. Every cell below comes from the
same fits, so nothing moves except the two axes.

\begin{table}[h]\centering
\caption{Net ARR. Flat fee is $\kappa_0$ alone; full adds slippage, impact and
borrow.}
\adjustbox{max width=\textwidth}{\MatrixNetARR}
\end{table}

\begin{table}[h]\centering
\caption{Book turnover. Top-$k$ is structurally pinned at $n_{\mathrm{drop}}/k$;
mean-variance spends its budget.}
\adjustbox{max width=\textwidth}{\MatrixTurnover}
\end{table}

\paragraph{This reverses the construction recommendation.}
Under the full cost model Arm~B returns $+0.78\%$ with top-$k$ and $-0.60\%$
with mean-variance, and it is the \emph{only} configuration in the entire grid
that is profitable after full costs. The earlier preference for mean-variance
came from a sweep that compared mean-variance parameterisations against each
other and never against top-$k$; it answered which mean-variance, not whether
mean-variance.

The mechanism is visible in the turnover column. Top-$k$ trades $0.105$ per day
against mean-variance's $0.019$, and the same sweep found net IR rising
monotonically as the turnover budget is loosened ($+0.2619$ at $0.004$/day,
$+0.2963$ at $0.008$, $+0.3344$ at $0.02$). This alpha decays fast enough that
trading it is worth what trading it costs. Mean-variance is the more robust
construction --- the full cost model takes only $0.55$pp from the baseline under
it, against $4.93$pp under top-$k$ --- but robustness is not the objective when
the signal being protected is short-lived.

\paragraph{An unreconciled discrepancy.}
Our flat-fee top-$k$ cell puts Arm~B at $+6.28\%$, while Qlib's own backtest
puts the same library, the same $50/5$ construction and the same commission at
$+23.55\%$. Both are excess returns over CSI~300 net of cost, so they are
nominally the same quantity. The two engines differ in ways that could account
for it --- this one applies a prediction-scale calibration, A-share price-limit
and suspension masks, and point-in-time index membership --- but the size of the
gap has not been explained, and until it is, neither number should be published
as \emph{the} flat-fee result. The comparisons \emph{within} this grid are
unaffected: they all come from one engine and one set of fits, so they remain
valid for choosing between constructions.

\section{What this does not show}

\begin{itemize}\itemsep2pt\parskip0pt
  \item \textbf{The turnover budget binds for every row.} Under mean-variance
        turnover is $0.0192$--$0.0195$ and cost $0.61$--$0.65$~bps/day across all
        four configurations, because $g$ enforces the same cap on all of them;
        under top-$k$ it is pinned at $n_{\mathrm{drop}}/k$ for the same reason.
        Either way Arm~B's objective is supposed to find cheaper-to-trade
        factors and neither construction lets it express that. The advantage
        measured above is therefore \emph{in spite of} the construction, not
        because of it. Moving the cost inside the objective, so turnover becomes
        endogenous, is the change that would let the mechanism show itself.
  \item \textbf{One mining run per arm.} See the generation-noise caveat above.
  \item \textbf{Arm~B was mined before the unary-minus parser fix.} Every
        expression whose root was a negation counted zero nodes and was silently
        rejected, so Arm~B's library was drawn from a restricted proposal space.
        Whatever it achieves here, it achieved without access to a natural class
        of mutations.
  \item \textbf{Three of Arm~B's 156 admitted factors were never computable}
        and are absent from its design matrix; the count in the table is the
        number that priced the book, not the number in the library.
\end{itemize}

\end{document}
